\chapter{Edges, Features, and the 2D Convolution}

In the last two chapters we learned to read an image as a two-dimensional
signal --- a grid of brightness values $f[m,n]$ --- and we learned what its
2D Fourier transform says about it: low frequencies are the broad washes of
tone, high frequencies are the fine detail, the texture, and above all the
\emph{edges}. This chapter is about going after those edges directly, in the
spatial domain, with a single tool: the \textbf{2D convolution}. We will see
that much of classical computer vision --- finding boundaries, detecting
corners, blurring, sharpening, spotting blobs --- is just convolution with a
cleverly chosen little grid of numbers called a \emph{kernel}.

That local kernel operation also carries this chapter far beyond classic image
processing. The 2D convolution is the atomic operation of the convolutional neural
network (CNN), the architecture behind essentially all modern image
recognition. The difference between a 1980s edge detector and a 2020s vision
model is not the operation --- both convolve an image with small kernels. The
difference is \emph{where the kernels come from}: a human designed the old
ones from geometry; gradient descent \emph{learns} the new ones from data. By
the end of this chapter you will see that both rest on the same operation.

\begin{goalbox}
Our wrist tracker is about motion, not images --- so why study image edges?
Three reasons. First, the 2D convolution is the same flip-shift-multiply-sum
machine from Chapter~13, now run over a grid, and seeing it twice cements it.
Second, when we turn a window of IMU or EMG samples into a small 2D
``time $\times$ channel'' array (Chapter~38), the very same convolutional
layers that read images will read our gesture data. Third, the idea of a
\emph{learned filter bank} replacing hand-designed features is exactly how
modern gesture recognizers work. This chapter connects image processing to
gesture recognition.
\end{goalbox}

\chapterline

\section{An edge is a place where brightness changes fast}
\label{ch32:what-is-an-edge}

Look at any photograph and ask where the ``edges'' are. Your eye lands on the
boundary between an object and its background, between a shadow and the lit
floor, between the dark trunk of a tree and the bright sky. In every case the
defining feature is the same: as you walk across that boundary, the brightness
\emph{changes sharply over a short distance}. Inside a region of uniform tone,
brightness barely changes; at an edge, it jumps.

That observation suggests the strategy. ``Changes fast'' is the
language of the \emph{derivative}. In one dimension, a function changes fast
where its derivative is large in magnitude. An image is a function of two
variables, so the right object is the \textbf{gradient} --- the pair of partial
derivatives,
\begin{equation}
\label{ch32:gradient-cont}
\nabla f(x,y) \;=\; \left( \frac{\partial f}{\partial x},\; \frac{\partial f}{\partial y} \right).
\end{equation}
The gradient is a vector. Its magnitude $\norm{\nabla f}$ tells you \emph{how
fast} brightness is changing at that point, and its direction tells you
\emph{which way} the brightness increases most steeply --- which is always
perpendicular to the edge itself.

\begin{intuitionbox}
Think of the image as an elevation map: bright pixels are hilltops, dark pixels are
valleys. The gradient at a point is the direction a marble would roll if you
flipped the surface upside down --- straight uphill --- and its magnitude is
how steep that slope is. A flat plateau (uniform region) has zero gradient. A
cliff (an edge) has a huge gradient pointing over the edge. To find edges, find
the cliffs: find where $\norm{\nabla f}$ is large.
\end{intuitionbox}

\begin{figure}[H]
\centering
\begin{tikzpicture}
\begin{axis}[
  width=12.5cm, height=5.2cm,
  xlabel={position $x$ across a row of pixels}, ylabel={brightness},
  xmin=0, xmax=10, ymin=-1.4, ymax=2.4,
  xtick={0,2,4,6,8,10}, ytick={0,1,2},
  axis lines=left, ymajorgrids, grid style={gray!18},
  legend style={at={(0.02,0.97)},anchor=north west,font=\small,draw=gray!40},
  legend cell align=left,
]
% a smoothed step edge: f(x)
\addplot[thick, color=sigblue, samples=160, domain=0:10]
  {1/(1+exp(-3*(x-5)))};
\addlegendentry{$f$ (brightness)}
% its derivative: a bump at the edge
\addplot[thick, color=specred, samples=160, domain=0:10]
  {3*exp(-3*(x-5))/(1+exp(-3*(x-5)))^2};
\addlegendentry{$f'$ (gradient magnitude)}
\addplot[dashed, gray, samples=2, domain=0:10]{0};
\end{axis}
\end{tikzpicture}
\caption{One row of pixels crossing an edge. The brightness $f$ rises from dark
to light over a short stretch (a smoothed step). Its derivative $f'$ is near
zero on the flat parts and \emph{peaks exactly at the edge}. Detecting edges
means finding the peaks of the derivative.}
\label{ch32:fig-step-edge}
\end{figure}

\subsection{Discrete derivatives: differences between neighbors}

Our image is sampled --- a grid of integers, not a smooth surface --- so we
cannot take a continuous derivative. Instead we go straight back to the definition:
a derivative is a limit of finite differences, and on a pixel grid of unit spacing we keep the
difference. The cleanest estimate of $\partial f/\partial x$ at pixel $[m,n]$
is the \textbf{central difference},
\begin{equation}
\label{ch32:central-diff}
\frac{\partial f}{\partial x}\bigg|_{[m,n]} \;\approx\; \frac{f[m,n+1] - f[m,n-1]}{2},
\end{equation}
the brightness of the pixel to the right minus the pixel to the left (the $2$
in the denominator accounts for the two-pixel span; in practice it is often
dropped because we only care about \emph{relative} edge strength). The vertical
derivative is the same idea down a column.

Notice how Equation~\eqref{ch32:central-diff} can be expressed directly as a
compact convolution. Lay the coefficients $\tfrac12(\,+1,\,0,\,-1\,)$ over the
three pixels $(n+1, n, n-1)$, multiply, and sum. That little row of numbers is a
\emph{kernel}. Taking a derivative \emph{is} convolving with a kernel. Once we
see that, every operation in this chapter becomes ``pick a kernel and
convolve.''

\section{The 2D convolution, concretely}
\label{ch32:2d-conv}

Recall the 1D convolution from Chapter~13: flip one signal, slide it across the
other, and at each shift multiply overlapping samples and add them up. The 2D
version is identical in spirit --- we just slide a small grid over a big grid.

\begin{definition}[2D convolution]
\label{ch32:def-conv}
Given an image $f[m,n]$ and a kernel $h[i,j]$, their convolution is
\begin{equation}
\label{ch32:conv2d}
(f \conv h)[m,n] \;=\; \sum_{i}\sum_{j} f[m-i,\,n-j]\,h[i,j].
\end{equation}
The output pixel at $[m,n]$ is a weighted sum of the input pixels in a
neighborhood, with the weights given by the (flipped) kernel.
\end{definition}

\begin{happeningbox}
Here is what the formula does. Park the
kernel --- a small grid, usually $3\times3$ or $5\times5$ --- on top of one
pixel of the image, centered there. Multiply each kernel weight by the image
pixel beneath it. Add up all those products. That single number is the output
pixel. Now slide the kernel one step over and repeat, for every pixel in the
image. The output is a new image the same size as the input. Every output
pixel is a weighted sum of its neighborhood, with the weights supplied by the
kernel.
\end{happeningbox}

\begin{figure}[H]
\centering
\begin{tikzpicture}[scale=0.58, >=Latex, font=\small]
  % Stage 1: Input image grid with sliding window
  \begin{scope}[shift={(0,0)}]
    \node[above, font=\small\bfseries] at (2.5,5.3) {Input Image $f[m,n]$};
    
    % Draw 5x5 grid
    \foreach \x in {0,...,4}{
      \foreach \y in {0,...,4}{
        \fill[gray!6] (\x,\y) rectangle (\x+1,\y+1);
        \draw[gray!35] (\x,\y) rectangle (\x+1,\y+1);
      }
    }
    
    % Highlight 3x3 receptive field patch at (1,1) to (4,4)
    \fill[sigblue!20] (1,1) rectangle (4,4);
    \draw[sigblue, line width=1.5pt] (1,1) rectangle (4,4);
    
    % Sample values inside the patch
    \node[font=\footnotesize] at (1.5,3.5) {$10$};
    \node[font=\footnotesize] at (2.5,3.5) {$10$};
    \node[font=\footnotesize] at (3.5,3.5) {$90$};
    \node[font=\footnotesize] at (1.5,2.5) {$10$};
    \node[font=\footnotesize\bfseries, sigblue!90!black] at (2.5,2.5) {$10$};
    \node[font=\footnotesize] at (3.5,2.5) {$90$};
    \node[font=\footnotesize] at (1.5,1.5) {$10$};
    \node[font=\footnotesize] at (2.5,1.5) {$10$};
    \node[font=\footnotesize] at (3.5,1.5) {$90$};

    % Sliding motion arrow
    \draw[->, thick, sigblue] (2.2, 4.3) -- (3.8, 4.3)
      node[midway, above=1pt, font=\tiny\bfseries, sigblue] {slide $\to$};
    \draw[->, thick, dashed, sigblue!70] (4.3, 2.8) -- (4.3, 1.2);

    \node[below, font=\footnotesize, align=center] at (2.5,-0.3) {
      $3\times3$ Receptive Field\\centered at pixel $[m,n]$
    };
  \end{scope}

  % Transition arrow: elementwise multiply
  \draw[->, thick, gray!70] (5.8, 2.5) -- (8.2, 2.5)
    node[midway, above=3pt, font=\footnotesize\bfseries, text=black] {$\odot$ overlay}
    node[midway, below=3pt, font=\scriptsize, text=gray!80!black] {multiply};

  % Stage 2: 3x3 Kernel
  \begin{scope}[shift={(8.9,1.0)}]
    \node[above, font=\small\bfseries] at (1.5,3.3) {Kernel $h[i,j]$};
    
    \foreach \x in {0,1,2}{
      \foreach \y in {0,1,2}{
        \fill[specred!12] (\x,\y) rectangle (\x+1,\y+1);
        \draw[specred, thick] (\x,\y) rectangle (\x+1,\y+1);
      }
    }
    % Kernel weights (horizontal Sobel / difference)
    \node[font=\footnotesize] at (0.5,2.5) {$-1$};
    \node[font=\footnotesize] at (1.5,2.5) {$0$};
    \node[font=\footnotesize] at (2.5,2.5) {$+1$};
    \node[font=\footnotesize] at (0.5,1.5) {$-2$};
    \node[font=\footnotesize\bfseries, specred!90!black] at (1.5,1.5) {$0$};
    \node[font=\footnotesize] at (2.5,1.5) {$+2$};
    \node[font=\footnotesize] at (0.5,0.5) {$-1$};
    \node[font=\footnotesize] at (1.5,0.5) {$0$};
    \node[font=\footnotesize] at (2.5,0.5) {$+1$};

    \node[below, font=\footnotesize, align=center] at (1.5,-0.3) {
      Fixed weights\\sum of products $= \mathbf{320}$
    };
  \end{scope}

  % Transition arrow: sum
  \draw[->, thick, gray!70] (12.6, 2.5) -- (15.0, 2.5)
    node[midway, above=3pt, font=\footnotesize\bfseries, text=black] {$\sum$}
    node[midway, below=3pt, font=\scriptsize, text=gray!80!black] {accumulate};

  % Stage 3: Output feature map
  \begin{scope}[shift={(15.7,0)}]
    \node[above, font=\small\bfseries] at (2.5,5.3) {Output Feature Map $g$};
    
    % Draw 5x5 grid
    \foreach \x in {0,...,4}{
      \foreach \y in {0,...,4}{
        \fill[gray!6] (\x,\y) rectangle (\x+1,\y+1);
        \draw[gray!35] (\x,\y) rectangle (\x+1,\y+1);
      }
    }
    
    % Previously scanned pixels (dim green)
    \foreach \x in {1,...,3}{
      \foreach \y in {3}{
        \fill[targetgreen!12] (\x,\y) rectangle (\x+1,\y+1);
      }
    }
    \fill[targetgreen!12] (1,2) rectangle (2,3);

    % Current output pixel at (2,2)
    \fill[targetgreen!40] (2,2) rectangle (3,3);
    \draw[targetgreen!90!black, line width=1.5pt] (2,2) rectangle (3,3);
    \node[font=\footnotesize\bfseries, targetgreen!90!black] at (2.5,2.5) {$320$};

    \node[below, font=\footnotesize, align=center] at (2.5,-0.3) {
      Single output pixel $g[m,n]$\\
      Sliding fills full feature map
    };
  \end{scope}
\end{tikzpicture}
\caption{The two-dimensional convolution operating as a sliding window across the pixel grid. At each location $[m,n]$, the $3\times3$ kernel is overlaid on the local neighborhood of the input image $f$, multiplied elementwise with the underlying pixel values, and accumulated into a single scalar in the output feature map $g[m,n]$. Advancing the window across columns and rows constructs the complete transformed image.}
\label{ch32:fig-conv-overlay}
\end{figure}

\begin{trapbox}
\textbf{Convolution flips the kernel; correlation does not.} The minus signs in
Eq.~\eqref{ch32:conv2d} flip the kernel both horizontally and vertically before
overlaying. Most image-processing code and \emph{all} deep-learning frameworks
actually implement \emph{cross-correlation},
$\sum_{i,j} f[m+i,n+j]\,h[i,j]$, which skips the flip. For a symmetric kernel
(a blur) the two are identical. For an antisymmetric kernel (a derivative) the
flip just negates the result --- a sign you can absorb into the kernel. So in
practice people say ``convolution'' and mean ``correlation,'' and it almost
never bites you. We will follow that convention but keep the distinction in
mind; it matters when you compare a textbook kernel to a library's output.
\end{trapbox}

\begin{whybox}
Why is convolution the \emph{right} operation, and not some other neighborhood
combination? Because it is the only operation that is both \emph{linear} (the
response to a sum of images is the sum of the responses) and
\emph{shift-invariant} (a feature looks the same wherever it appears in the
frame). An edge in the top-left corner and the same edge in the bottom-right
must produce the same response --- a face is a face wherever it sits in the
photo. Exactly as in Chapter~12, ``linear + shift-invariant'' forces the
operation to be convolution with some kernel. The kernel \emph{is} the system's
impulse response, now a little 2D patch.
\end{whybox}

\subsection{Boundaries, separability, and cost}

Two practical points before we start choosing kernels.

\textbf{Edges of the image.} When the kernel hangs off the side of the image,
some of its weights have no pixel beneath them. The standard fixes are to
\emph{zero-pad} (pretend the outside is black), \emph{replicate} the border
pixel outward, or \emph{reflect} the image across its edge. Reflection usually
looks best because it avoids inventing a fake dark border that would itself
register as an edge.

\textbf{Separability.} A kernel is \emph{separable} if it factors as an outer
product of a column and a row, $h[i,j] = c[i]\,r[j]$. Then a 2D convolution
splits into two cheap 1D convolutions: convolve every row with $r$, then every
column with $c$. This drops the cost per output pixel from $K^2$ multiplies to
$2K$ --- a large saving for big kernels, and the reason the Gaussian blur (which
is separable) is so fast. We will see that several of our edge kernels are
separable too.

\begin{example}[A $3\times3$ convolution by hand]
\label{ch32:ex-byhand}
Take a tiny image patch and the horizontal-difference kernel:
\[
f = \begin{pmatrix} 10 & 10 & 90 \\ 10 & 10 & 90 \\ 10 & 10 & 90 \end{pmatrix},
\qquad
h = \begin{pmatrix} -1 & 0 & 1 \\ -1 & 0 & 1 \\ -1 & 0 & 1 \end{pmatrix}.
\]
Centered on the middle pixel, cross-correlation gives
\[
\sum h_{ij} f_{ij}
= (-1)(10) + (0)(10) + (1)(90) + (-1)(10) + 0 + (1)(90) + (-1)(10) + 0 + (1)(90)
= 240.
\]
A large positive number --- this kernel has detected the strong dark-to-light
edge running down the right side of the patch. Slide it to a flat patch where
all pixels equal $10$ and the sum is $0$: no edge, no response. The kernel
\emph{is} an edge detector, and the arithmetic is just a weighted sum of nine
numbers.
\end{example}

\section{From differences to the Sobel operator}
\label{ch32:sobel}

The crude central-difference kernel $(\,+1,\,0,\,-1\,)$ works, but it has a
flaw: it is brutally sensitive to noise. A single noisy pixel produces a fake
edge. The cure is to \emph{average in the direction along the edge while
differencing across it}. That is exactly what the famous \textbf{Sobel
operator} does. Its horizontal-gradient kernel is
\begin{equation}
\label{ch32:sobel-x}
G_x \;=\;
\begin{pmatrix} -1 & 0 & +1 \\ -2 & 0 & +2 \\ -1 & 0 & +1 \end{pmatrix},
\qquad
G_y \;=\;
\begin{pmatrix} -1 & -2 & -1 \\ \phantom{-}0 & \phantom{-}0 & \phantom{-}0 \\ +1 & +2 & +1 \end{pmatrix}.
\end{equation}

Read $G_x$ as a product of two intentions. \emph{Across} the kernel (left to
right) it is the difference $(-1,0,+1)$ --- a derivative. \emph{Down} the kernel
it is the smoothing weights $(1,2,1)$ --- a little blur. In fact $G_x$ factors
exactly:
\begin{equation}
\label{ch32:sobel-sep}
G_x =
\begin{pmatrix} 1 \\ 2 \\ 1 \end{pmatrix}
\begin{pmatrix} -1 & 0 & +1 \end{pmatrix},
\end{equation}
so it is separable: smooth vertically with $(1,2,1)$, then difference
horizontally with $(-1,0,+1)$. \emph{Differentiate across the edge, blur along
it}: that is the basic design principle of a good gradient operator.

\begin{whybox}
Why blend blurring with differencing? Differentiation \emph{amplifies} high
frequencies --- and noise is the highest-frequency content of all (Chapter~20).
A pure difference kernel is a high-pass filter, so it boosts noise enormously.
Blurring is low-pass; it suppresses noise. Sobel multiplies a low-pass along
one axis by a high-pass along the other, giving a \emph{band-pass} response
that is tuned to real edges (which are spatially extended structures) while
rejecting isolated noisy pixels (which are not). This is the same
smooth-then-differentiate logic behind the derivative-of-Gaussian filter we
meet in the next section.
\end{whybox}

Once we have $G_x$ and $G_y$ as two convolved images, we combine them into the
gradient magnitude and direction at every pixel:
\begin{equation}
\label{ch32:grad-mag}
\norm{\nabla f} = \sqrt{G_x^2 + G_y^2}, \qquad
\theta = \operatorname{atan2}(G_y,\, G_x).
\end{equation}
The magnitude image is bright exactly along edges and near-black everywhere
else; the direction tells us each edge's orientation. Thresholding the
magnitude gives a first, rough edge map.

\begin{figure}[H]
\centering
\begin{tikzpicture}
\begin{axis}[
  width=12.5cm, height=5cm,
  xlabel={position $x$ across the row}, ylabel={value},
  xmin=0, xmax=12, ymin=-1.7, ymax=2.6,
  xtick={0,2,4,6,8,10,12}, ytick={-1,0,1,2},
  axis lines=left, ymajorgrids, grid style={gray!18},
  legend style={at={(0.98,0.97)},anchor=north east,font=\small,draw=gray!40,fill=white,inner sep=4pt},
  legend cell align=left,
]
% noisy step
\addplot[thick, color=sigblue, samples=200, domain=0:12]
  {1/(1+exp(-2.5*(x-6))) + 0.12*sin(deg(11*x)) + 0.08*cos(deg(23*x))};
\addlegendentry{noisy brightness}
% naive derivative (amplifies noise) -- emulate with a wiggly bump
\addplot[thin, color=gray, samples=200, domain=0:12]
  {2.5*exp(-2.5*(x-6))/(1+exp(-2.5*(x-6)))^2 + 0.35*cos(deg(23*x))};
\addlegendentry{raw difference}
% smoothed derivative (Sobel-like) -- clean bump
\addplot[thick, color=specred, samples=200, domain=0:12]
  {2.5*exp(-2.5*(x-6))/(1+exp(-2.5*(x-6)))^2};
\addlegendentry{Sobel (smoothed)}
\addplot[dashed, gray, samples=2, domain=0:12]{0};
\end{axis}
\end{tikzpicture}
\caption{Why Sobel smooths along the edge. The raw difference of a noisy
brightness profile (gray) is itself noisy --- full of spurious bumps that would
become fake edges. Sobel's built-in $(1,2,1)$ smoothing yields a clean bump
(red) located right at the true edge.}
\label{ch32:fig-sobel-smooth}
\end{figure}

\begin{lstlisting}[language=Python,caption={Sobel edge magnitude in NumPy/SciPy.},label={ch32:lst-sobel}]
import numpy as np
from scipy.ndimage import convolve

Gx = np.array([[-1, 0, 1],
               [-2, 0, 2],
               [-1, 0, 1]], dtype=float)
Gy = Gx.T  # transpose gives the vertical kernel

def sobel_edges(img):
    img = img.astype(float)
    gx = convolve(img, Gx, mode="reflect")   # reflect padding at borders
    gy = convolve(img, Gy, mode="reflect")
    mag = np.hypot(gx, gy)                    # sqrt(gx**2 + gy**2)
    return mag / mag.max()                    # normalize to [0, 1]
\end{lstlisting}

\begin{applicationbox}
Sobel and its cousins (Prewitt, with equal $(1,1,1)$ smoothing; Scharr, with
optimized $(3,10,3)$ weights for better rotational accuracy) are everywhere in
real systems: lane-line detection in cars, defect inspection on assembly lines,
fingerprint ridge enhancement, and the first preprocessing step of countless
medical-imaging pipelines. They run in real time on the cheapest hardware
because they are just nine multiplies per pixel --- which is precisely why they
survived into the deep-learning era as the kind of filter a network's first
layer reliably re-discovers on its own.
\end{applicationbox}

\section{Canny: turning a gradient into clean edges}
\label{ch32:canny}

A thresholded Sobel magnitude is messy: edges come out as thick smears, faint
edges flicker in and out depending on the threshold, and noise speckles the
result. In 1986 John Canny asked what an \emph{optimal} edge detector would
look like and derived a multi-stage procedure that is still the default in
practice. We will not redo his variational derivation, but the four stages are
intuitive once you have the gradient picture (summarized in Figure~\ref{ch32:fig-canny-pipeline}).

\begin{enumerate}
\item \textbf{Smooth with a Gaussian.} Convolve the image with a Gaussian blur
first, to tame noise before differentiating. The blur width $\sigma$ sets the
scale of edges you care about: small $\sigma$ finds fine detail, large $\sigma$
finds only the big boundaries.

\item \textbf{Compute the gradient.} Apply Sobel to get $\norm{\nabla f}$ and
$\theta$ at every pixel. (Steps 1 and 2 together are a
\emph{derivative-of-Gaussian} filter --- one convolution, by associativity.)

\item \textbf{Non-maximum suppression.} An edge should be one pixel thick, not
a fat ridge. Walk along the gradient direction; keep a pixel only if its
magnitude is a \emph{local maximum} compared to its two neighbors across the
edge. This thins the smear down to a crisp line at the ridge crest.

\item \textbf{Hysteresis thresholding.} Use \emph{two} thresholds. Pixels above
the high threshold are definitely edges. Pixels below the low threshold are
definitely not. Pixels in between count as edges only if they connect to a
strong edge. This links up broken contours while rejecting isolated noise ---
the trick that makes Canny output look like clean line drawings.
\end{enumerate}

\begin{figure}[H]
\centering
\begin{tikzpicture}[
  stage/.style={rounded corners=4pt, draw, thick, minimum width=2.3cm, minimum height=1.35cm, align=center, font=\small},
  io/.style={rounded corners=3pt, draw=gray!60, fill=gray!8, minimum width=1.45cm, minimum height=1.2cm, align=center, font=\footnotesize},
  arr/.style={-{Latex[length=2.2mm]}, thick, color=gray!60},
  node distance=4.5mm
]
\node[io] (in) {Input\\image $f$};
\node[stage, fill=sourceblue!10, draw=sourceblue, right=of in] (s1) {\textbf{1. Gaussian}\\\scriptsize noise smoothing\\\scriptsize scale $\sigma$};
\node[stage, fill=quotientorange!10, draw=quotientorange, right=of s1] (s2) {\textbf{2. Gradient}\\\scriptsize Sobel filter\\\scriptsize $\norm{\nabla f},\, \theta$};
\node[stage, fill=structurepurple!10, draw=structurepurple, right=of s2] (s3) {\textbf{3. Non-max}\\\scriptsize suppress smear\\\scriptsize 1-pixel crest};
\node[stage, fill=targetgreen!10, draw=targetgreen, right=of s3] (s4) {\textbf{4. Hysteresis}\\\scriptsize dual threshold\\\scriptsize $T_{\text{lo}},\, T_{\text{hi}}$};
\node[io, right=of s4] (out) {Clean\\edges};

\draw[arr] (in) -- (s1);
\draw[arr] (s1) -- (s2);
\draw[arr] (s2) -- (s3);
\draw[arr] (s3) -- (s4);
\draw[arr] (s4) -- (out);
\end{tikzpicture}
\caption{The Canny edge detection pipeline. Raw image data is smoothed to tame noise at scale $\sigma$, differentiated to find gradient vector magnitude and orientation, thinned to single-pixel ridge crests via non-maximum suppression, and linked into continuous contours using dual-threshold hysteresis.}
\label{ch32:fig-canny-pipeline}
\end{figure}

\begin{intuitionbox}
Sobel answers ``how steep is the cliff here?'' Canny answers ``draw me the exact
ridgeline of the cliffs, one pixel wide, with no gaps and no false starts.''
The first is a measurement; the second is a decision. Non-maximum suppression
is the ``thin it'' decision; hysteresis is the ``connect it'' decision. Both are
heuristics layered on top of the same gradient convolution.
\end{intuitionbox}

\begin{connectionbox}
The derivative-of-Gaussian filter from step 2 is a clean application of
calculus and linear algebra (Chapter~42). By the convolution theorem
(Chapter~11), \emph{differentiating a blurred image equals convolving with the
derivative of the blur}: $\frac{\dd}{\dd x}(g \conv f) = (\frac{\dd g}{\dd x})
\conv f$. So we never differentiate the noisy image at all --- we differentiate
the smooth Gaussian once, analytically, and convolve that fixed kernel with the
image. Calculus does the hard part on paper; convolution does the cheap part on
pixels.
\end{connectionbox}

\begin{historybox}
The search for an ``optimal'' edge detector took about two decades and drew on
both neuroscience and engineering.
In 1968, Irwin Sobel and Gary Feldman at the Stanford Artificial Intelligence
Project (SAIL) were trying to teach a robotic arm to segment simple wooden blocks
under harsh lab lights. The naive approach --- computing raw differences between
adjacent pixels --- produced large noise spikes, mistaking sensor grain for object
boundaries. Sobel and Feldman realized that differencing
must be balanced by transverse smoothing: differencing across the edge while
averaging along it, yielding the famous $3\times3$ separable operator.

A decade later at MIT, David Marr and Ellen Hildreth (1980) took inspiration from
David Hubel and Torsten Wiesel's Nobel-winning 1959 discovery that neurons in the
mammalian primary visual cortex act as oriented spatial filters. Marr and Hildreth
proposed finding the zero crossings of a Laplacian-of-Gaussian (LoG) filter. Yet
their method had its own problem: zero crossings
always form closed loops, forcing the algorithm to invent phantom contours even
in smooth gradients where no physical surface boundary existed.

The next major step came in 1986, when John Canny, also at MIT, recast
edge detection as a formal variational optimization problem. Instead of guessing
a stencil, Canny laid down three mathematical criteria:
low error rate (maximizing signal-to-noise ratio), high localization (the detected
ridge must lie as close as possible to the true physical center), and a single response
per edge (no split or multi-peaked duplicate contours). Solving the calculus of variations, Canny
showed that the first derivative of a Gaussian approximates this optimal filter to
within a fraction of a percent. To cure the broken-contour defect of single
thresholds, he introduced non-maximum suppression along the gradient normal and
dual-threshold hysteresis tracking.

For thirty years, Canny's detector, the Harris corner detector (1988), and
scale-space blob filters were the standard front end of computer vision, industrial
robotics, and satellite tracking. Hand-engineered convolution eventually reached its
limits: while translation-equivariant, fixed spatial kernels
cannot adapt to non-rigid deformations, illumination shifts, or 3D viewpoint changes.
In 1989, Yann LeCun at Bell Labs closed the loop by replacing hand-crafted stencils
with learned kernels optimized via backpropagation, producing the convolutional
neural network. As we move from still images to transform compression in Chapter~33
and wearable motion sensors from Chapter~34 onward, the same convolution operation
is reused as a temporal feature detector for real-time motion and gesture
recognition.
\end{historybox}

\section{Blurring versus differentiating: two faces of one operation}
\label{ch32:blur-vs-diff}

Two operations run through everything above. Blurring and differentiating are
both convolutions, but with opposite-looking kernels,
and they do opposite things in the frequency domain.

A \textbf{blur} kernel has all-positive weights that sum to $1$. The simplest
is the $3\times3$ box average (every weight $\tfrac19$); the better-behaved one
is the Gaussian,
\begin{equation}
\label{ch32:gauss-kernel}
g[i,j] \;\propto\; \exp\!\left(-\frac{i^2 + j^2}{2\sigma^2}\right),
\end{equation}
normalized so the weights sum to $1$. Because the weights sum to one, a blur
preserves the overall brightness and \emph{attenuates high frequencies} --- it
is a low-pass filter. The output is a smoothed, softened image.

\begin{figure}[H]
\centering
\begin{tikzpicture}
\begin{axis}[
  width=11cm, height=4.6cm,
  xlabel={offset from center (pixels)}, ylabel={weight},
  xmin=-4, xmax=4, ymin=0, ymax=0.62,
  xtick={-4,-2,0,2,4}, ytick={0,0.2,0.4},
  axis lines=left, ymajorgrids, grid style={gray!18},
  legend style={at={(0.98,0.97)},anchor=north east,font=\small,draw=gray!40,fill=white,inner sep=4pt},
  legend cell align=left,
]
% Gaussian 1D profile, sigma=1
\addplot[thick, color=sigblue, samples=120, domain=-4:4]
  {0.3989*exp(-x^2/2)};
\addlegendentry{Gaussian ($\sigma=1$)}
% box average over 5 taps (each 0.2) shown as a step
\addplot[thick, color=quotientorange, const plot, samples at={-4,-3,-2.5,-1.5,-0.5,0.5,1.5,2.5,3,4}]
  {0.2*(abs(x)<2.5)};
\addlegendentry{box average (5 taps)}
\end{axis}
\end{tikzpicture}
\caption{Two smoothing kernels in 1D. The box average gives every neighbor an
equal vote (a hard window); the Gaussian weights nearby pixels more than distant
ones (a soft window). Both sum to $1$, so both preserve brightness, but the
Gaussian's smooth taper avoids the ringing artifacts a sharp box edge produces.}
\label{ch32:fig-blur-kernels}
\end{figure}

A \textbf{derivative} kernel has weights that sum to \emph{zero} (look back at
Sobel: $-1-2-1+1+2+1 = 0$). A zero-sum kernel ignores constant brightness
entirely --- convolve it with a flat gray field and you get zero. It responds
only to \emph{changes}, which is why it is a high-pass filter and an edge
detector. Differentiating boosts high frequencies; blurring suppresses them.

\begin{focusbox}
The single most useful rule of thumb for reading any kernel:
\begin{itemize}
\item \textbf{Weights sum to $1$} $\Rightarrow$ it is a \emph{smoother /
low-pass} (blur, average). Preserves brightness, kills detail.
\item \textbf{Weights sum to $0$} $\Rightarrow$ it is a \emph{differentiator /
high-pass} (edge, derivative). Kills flat regions, keeps changes.
\item \textbf{Center positive, surround negative, summing to $1$}
$\Rightarrow$ it is a \emph{sharpener} --- original plus a dash of high-pass.
\end{itemize}
Glance at the numbers, add them up, and you already know what a kernel does
before you ever run it.
\end{focusbox}

The two combine to give the second-derivative edge detector, the
\textbf{Laplacian}, whose kernel is the discrete $\nabla^2 = \partial_{xx} +
\partial_{yy}$:
\begin{equation}
\label{ch32:laplacian}
L =
\begin{pmatrix} 0 & 1 & 0 \\ 1 & -4 & 1 \\ 0 & 1 & 0 \end{pmatrix}.
\end{equation}
It sums to zero (a differentiator), and instead of \emph{peaking} at an edge
like the first derivative, it \emph{crosses zero} there --- the second
derivative is positive on one side of the edge and negative on the other.
Finding ``zero crossings'' of the Laplacian is another classic edge detector
(Marr--Hildreth), and the Laplacian-of-Gaussian forms the core of the blob
detector we meet next.

\begin{figure}[H]
\centering
\begin{tikzpicture}
\begin{axis}[
  width=12.5cm, height=5cm,
  xlabel={position $x$}, ylabel={value},
  xmin=0, xmax=10, ymin=-2.4, ymax=2.6,
  xtick={0,2,4,6,8,10}, ytick={-2,-1,0,1,2},
  axis lines=left, ymajorgrids, grid style={gray!18},
  legend style={at={(0.02,0.03)},anchor=south west,font=\small,draw=gray!40},
  legend cell align=left,
]
\addplot[thick, color=sigblue, samples=160, domain=0:10]
  {1/(1+exp(-3*(x-5)))};
\addlegendentry{$f$ (edge)}
\addplot[thick, color=specred, samples=160, domain=0:10]
  {3*exp(-3*(x-5))/(1+exp(-3*(x-5)))^2};
\addlegendentry{$f'$ (peaks at edge)}
\addplot[thick, color=quotientorange, samples=200, domain=0:10]
  {9*exp(-3*(x-5))*(exp(-3*(x-5))-1)/(1+exp(-3*(x-5)))^3};
\addlegendentry{$f''$ (zero-crosses at edge)}
\addplot[dashed, gray, samples=2, domain=0:10]{0};
\end{axis}
\end{tikzpicture}
\caption{First versus second derivative at an edge. The first derivative $f'$
\emph{peaks} at the edge (find the maximum). The second derivative $f''$
\emph{crosses zero} at the edge (find the sign change). Sobel uses the first;
the Laplacian uses the second.}
\label{ch32:fig-first-second-deriv}
\end{figure}

\section{Beyond edges: corners and blobs}
\label{ch32:corners-blobs}

An edge is a 1D feature --- it has a direction, and you can slide along it
without anything changing. That makes a single edge a poor landmark: if you are
trying to track a point from frame to frame (which is exactly what motion and
gesture systems do), a point on an edge could be anywhere along that edge. What
you want is a feature that is \emph{distinctive in every direction} --- a
\textbf{corner}.

\subsection{Corners: the structure tensor}

The insight (Harris and Stephens, 1988; Moravec before them) is to ask: if I
nudge a little window in any direction, how much does the patch underneath
change? On a flat region, nudging changes nothing. Along an edge, nudging along
the edge changes nothing, but nudging across it changes a lot --- change in
\emph{one} direction. At a corner, nudging in \emph{any} direction changes the
patch a lot. Corners are where the brightness varies strongly along two
independent directions at once.

We make this precise with the \textbf{structure tensor}, built entirely from
the gradients we already have. Over a small window $W$ around a pixel, form
\begin{equation}
\label{ch32:structure-tensor}
M \;=\; \sum_{W}
\begin{pmatrix} G_x^2 & G_x G_y \\ G_x G_y & G_y^2 \end{pmatrix}.
\end{equation}
This $2\times2$ matrix summarizes how the gradient behaves in the window. Its
two eigenvalues $\lambda_1, \lambda_2$ classify the pixel:
\begin{itemize}
\item Both small $\Rightarrow$ flat region (no structure).
\item One large, one small $\Rightarrow$ an edge (strong variation in one
direction only).
\item Both large $\Rightarrow$ a corner (strong variation in two directions).
\end{itemize}
Computing eigenvalues per pixel is costly, so Harris uses the cheap
\emph{corner response} $R = \det M - k(\operatorname{tr} M)^2 = \lambda_1
\lambda_2 - k(\lambda_1 + \lambda_2)^2$, which is large and positive only when
both eigenvalues are large. Peaks of $R$ are corners.

\begin{connectionbox}
The structure tensor puts symmetric $2\times 2$ eigendecomposition (Chapter~42) directly to work.
The eigenvectors of $M$ point along the directions of least and most change;
the eigenvalues measure how much change. An edge is a near-singular $M$ (one
tiny eigenvalue); a corner is a well-conditioned $M$ (two healthy eigenvalues).
``Is this a corner?'' becomes ``is this $2\times2$ matrix
well-conditioned?'' --- and every ingredient ($G_x$, $G_y$, the window sum) is a
convolution.
\end{connectionbox}

\subsection{Blobs: the Laplacian of Gaussian across scale}

A \textbf{blob} is a roughly circular region that differs from its surroundings
--- a dot, a spot, a small object. The natural detector is the
\textbf{Laplacian-of-Gaussian} (LoG): blur with a Gaussian of width $\sigma$,
then take the Laplacian. The LoG kernel looks like a ``Mexican hat'': a positive
bump in the middle ringed by a negative moat. Convolved with the image, it
responds most strongly where a bright blob of size matched to $\sigma$ sits on a
darker background.

The key parameter is \emph{scale}. A blob has a size, and the LoG only responds
strongly when its $\sigma$ matches that size. So we run the LoG at a whole stack
of scales and look for peaks in both space \emph{and} scale --- the
$(x, y, \sigma)$ where the response is maximal tells us \emph{where} the blob is
and \emph{how big} it is. This scale-space search is the basis of classic
feature detectors like SIFT (which approximates the LoG by a cheap
\emph{difference of Gaussians}, two blurs subtracted).

\begin{figure}[H]
\centering
\begin{tikzpicture}
\begin{axis}[
  width=12cm, height=5cm,
  xlabel={position $x$ (1D slice through a blob)}, ylabel={LoG response},
  xmin=-6, xmax=6, ymin=-0.8, ymax=1.15,
  xtick={-6,-4,-2,0,2,4,6}, ytick={-0.5,0,0.5,1},
  axis lines=left, ymajorgrids, grid style={gray!18},
]
% Mexican-hat / Ricker wavelet: (1 - x^2) e^{-x^2/2}
\addplot[thick, color=structurepurple, samples=200, domain=-6:6]
  {(1 - x^2/2)*exp(-x^2/4)};
\addplot[dashed, gray, samples=2, domain=-6:6]{0};
\end{axis}
\end{tikzpicture}
\caption{A 1D slice of the Laplacian-of-Gaussian (``Mexican hat'') kernel: a
positive center flanked by negative lobes. It fires hardest on a bright blob
whose size matches the central lobe --- the basis of multi-scale blob
detection.}
\label{ch32:fig-log}
\end{figure}

\section{The 2D convolution as the core operation of a CNN}
\label{ch32:cnn}

Everything so far has been hand-designed: a human looked at the geometry of
edges, corners, and blobs, and wrote down kernels by reasoning. The key idea
of the convolutional neural network is to keep the
operation --- convolve the image with small kernels --- but \emph{stop choosing
the kernels by hand}. Instead, treat the kernel weights as unknowns and
\emph{learn} them from labeled examples by gradient descent.

\begin{happeningbox}
A convolutional layer is structurally identical to what we have been doing.
It holds a bank of small kernels (say sixty-four $3\times3$ kernels). It
convolves each one with the input image, producing sixty-four output images
called \emph{feature maps} --- each map is bright wherever its kernel's pattern
appears. A nonlinearity (usually ReLU: clamp negatives to zero) is applied,
and the maps are downsampled (\emph{pooling}). Then another convolutional layer
does the same to those feature maps, and another, stacking dozens deep. The
only thing the network ``knows'' at the start is that it will convolve; the
\emph{contents} of every kernel are learned.
\end{happeningbox}

The reason this works --- and the reason convolution rather than a fully
connected layer --- comes straight from the geometry we built earlier.

\begin{whybox}
Why convolution instead of a generic neural-network layer? Two priors baked
into the architecture, both of which we have already justified.
\begin{itemize}
\item \textbf{Locality.} A pixel's meaning is set by its neighbors, not by a
pixel on the far side of the image. A small kernel looks only at a local patch.
\item \textbf{Translation invariance / weight sharing.} The same kernel slides
over the whole image, so a feature is detected identically wherever it appears
--- exactly the shift-invariance argument from Section~\ref{ch32:2d-conv}. This
also slashes the number of parameters: a $3\times3$ kernel has nine weights, no
matter how big the image, versus millions for a fully connected layer.
\end{itemize}
These are the same two assumptions --- linear and shift-invariant --- that
forced us to convolution in the first place. The CNN simply lets the data fill
in the kernel weights that those assumptions left open.
\end{whybox}

\subsection{Learned filters rediscover the classics}

When you train a CNN on
natural images and then \emph{look at} the kernels its first convolutional
layer learned, you do not see random noise. You see oriented edge detectors at
various angles, color-contrast detectors, small blobs, and little textures ---
the learned first-layer filters look very much like Sobel kernels, Gabor
filters, and Laplacians of Gaussian. The network, given only a pile of photos
and a classification goal, \emph{rediscovers the very operators that Sobel,
Canny, Harris, and Marr designed by hand.}

\begin{aestheticbox}
This convergence has a simple cause: natural images are built from oriented edges and blobs at multiple
scales --- that is the statistical structure of photographs of a world made of
surfaces and objects under light. Any system that must efficiently represent
such images, whether a human visual cortex (whose first-stage cells Hubel and
Wiesel found to be oriented edge detectors in 1959), a hand-designed pipeline,
or a trained CNN, ends up with similar first-layer filters. Evolution,
engineering, and optimization arrive at the same filters because the statistics
of the images, not the designer, determine them.
\end{aestheticbox}

\begin{connectionbox}
In Chapter~39 we turn to how such model parameters are fit from labeled data:
backpropagation computes the gradient of a loss with respect to each kernel
weight, and gradient descent nudges them. The key fact that makes it work is
that \emph{convolution is differentiable in its weights} --- and the gradient of a convolution is itself a convolution. The forward pass, the
backward pass, and the weight update are all the same flip-shift-multiply-sum
machine we built in Chapter~13, run in three directions.
\end{connectionbox}

\begin{figure}[H]
\centering
\begin{tikzpicture}[node distance=4mm and 9mm, font=\small,
  every node/.style={align=center}]
\tikzset{
  blk/.style={draw, rounded corners, minimum height=10mm, minimum width=15mm,
    fill=gray!6, thick},
  conv/.style={blk, fill=sigblue!12, draw=sigblue!70},
  pool/.style={blk, fill=targetgreen!14, draw=targetgreen!70},
  fc/.style={blk, fill=quotientorange!14, draw=quotientorange!70},
}
\node[blk] (in) {image\\$H\times W$};
\node[conv, right=of in] (c1) {conv\\$+$ReLU};
\node[pool, right=of c1] (p1) {pool};
\node[conv, right=of p1] (c2) {conv\\$+$ReLU};
\node[pool, right=of c2] (p2) {pool};
\node[fc, right=of p2] (fc) {dense\\$\to$ class};
\foreach \a/\b in {in/c1, c1/p1, p1/c2, c2/p2, p2/fc}
  \draw[-{Stealth},thick,gray] (\a) -- (\b);
\node[font=\scriptsize,gray,below=1mm of c1] {learned\\edges};
\node[font=\scriptsize,gray,below=1mm of c2] {learned\\parts};
\end{tikzpicture}
\caption{A small CNN as a stack of convolutions. Each \texttt{conv} block is a
bank of learned 2D kernels --- the same operation as a Sobel filter, but with
weights set by training. Early layers learn edge-like filters; deeper layers
combine them into parts and objects; a final dense layer reads out a class.}
\label{ch32:fig-cnn}
\end{figure}

\section{Carrying the idea to the wrist}
\label{ch32:to-wrist}

\begin{goalbox}
The wrist tracker never sees an image --- so where does the 2D convolution land?
In two concrete places.

\textbf{1. Time $\times$ channel maps.} Take a one-second window of our sensor
data and stack the channels into a 2D array: rows are channels (3 accelerometer
axes, 3 gyroscope axes, several EMG electrodes), columns are time samples. That
array \emph{is} a small image, and a CNN's 2D kernels can sweep across it,
learning filters that detect, say, ``a sharp flexion burst on EMG channel~2
coincident with a downward accelerometer spike.'' This is exactly how many
published gesture and human-activity-recognition models work, and it is the
same convolution from this chapter.

\textbf{2. 1D convolution as the natural form.} More often, each sensor channel
is convolved on its own with a small 1D kernel (a \emph{1D CNN}). That is
the 1D special case of Definition~\ref{ch32:def-conv}: the same
learned-filter idea, sliding a tiny window along the time axis to detect
temporal motifs --- the onset of a tap, the plateau of a held pose.

In both cases the idea from this chapter carries over unchanged: a small, shared,
learned kernel sliding over the signal, detecting a local pattern wherever it
occurs.
\end{goalbox}

\begin{trapbox}
\textbf{Know what convolution can and cannot recover.} A learned
convolutional filter is only as informative as the signal it sees. No amount of
clever kernels on a \emph{wrist-worn IMU alone} will reconstruct continuous
individual finger-joint angles --- the wrist simply does not carry that
information, and a CNN cannot conjure it from data that is not there. What 2D/1D
convolutional features genuinely deliver from the wrist is robust detection of
\emph{coarse, discrete} gestures (a flick, a twist, a tap, a hold) and, with
forearm \emph{EMG or force myography}, discrimination among a modest set of
finger and grip patterns. Continuous finger \emph{pose} needs sensors on the
fingers themselves (instrumented rings or gloves). The convolution is a
pattern matcher, not a source of information that the sensor never captured.
\end{trapbox}

\begin{buildbox}
A practical first-layer feature extractor for the wrist, runnable on a Pi or a
beefy microcontroller, mirrors this chapter exactly. Treat a window of IMU
samples as a $C \times T$ array (channels by time) and convolve with a handful
of small learned 1D kernels, then pool. The skeleton:
\begin{lstlisting}[language=Python,caption={A 1D conv feature stage for IMU windows (PyTorch sketch).},label={ch32:lst-imu-cnn}]
import torch.nn as nn

# input: (batch, C_channels=6, T_samples=128)  -- a window of accel+gyro
feature = nn.Sequential(
    nn.Conv1d(6, 16, kernel_size=7, padding=3),  # 16 learned temporal kernels
    nn.ReLU(),
    nn.MaxPool1d(2),                             # downsample in time
    nn.Conv1d(16, 32, kernel_size=5, padding=2), # combine into motifs
    nn.ReLU(),
    nn.AdaptiveAvgPool1d(1),                     # one vector per window
)
# feed the resulting vector to a small classifier over your gesture set
\end{lstlisting}
Each \texttt{Conv1d} kernel is the 1D sibling of a Sobel row. Train it on
labeled gestures and inspect the learned kernels: you will find some that look
like edge detectors (firing at motion \emph{onsets}) and some that look like
smoothers (averaging out tremor) --- the same two faces of convolution from
Section~\ref{ch32:blur-vs-diff}.
\end{buildbox}

\section{Summary}
\label{ch32:summary}

\begin{itemize}
\item An \textbf{edge} is where brightness changes fast; the \textbf{gradient}
$\nabla f$ measures that change in magnitude and direction.
\item Estimating derivatives on a grid is a small \textbf{convolution}: a
derivative is just a kernel of differences.
\item The \textbf{2D convolution} (Def.~\ref{ch32:def-conv}) is overlay-
multiply-sum of a small kernel over the image --- the unique linear,
shift-invariant neighborhood operation.
\item \textbf{Sobel} differentiates across an edge while smoothing along it
(separable, noise-robust); \textbf{Canny} thins and links the gradient into
clean one-pixel contours.
\item A kernel summing to \textbf{1} blurs (low-pass); summing to \textbf{0}
differentiates (high-pass). The \textbf{Laplacian} is a second-derivative
detector whose zero crossings mark edges.
\item \textbf{Corners} (structure-tensor eigenvalues both large) and
\textbf{blobs} (Laplacian-of-Gaussian peaks across scale) are more distinctive
features, all built from gradient convolutions.
\item A \textbf{CNN} keeps the convolution but \emph{learns} the kernels; its
first layer reliably rediscovers edge, blob, and texture detectors --- the same
operators engineers designed by hand.
\item For the wrist, the same idea appears as \textbf{1D/2D convolutions over
time $\times$ channel windows} --- a reliable detector of coarse gestures, but
not a substitute for sensors that the wrist alone cannot replace.
\item \textbf{Looking forward:} While convolution extracts localized features like
edges and spots, Chapter~33 (Section~\ref{ch33:why-transform}) works on whole image blocks,
rotating pixel blocks into orthonormal cosine basis coordinates to compress
photographs down to a fraction of their raw storage.
\end{itemize}

\section{Exercises}
\label{ch32:exercises}

\begin{exercise}
\label{ch32:ex-byhand2}
Take the $3\times3$ patch $f = \big(\begin{smallmatrix} 5 & 5 & 5 \\ 5 & 5 & 5
\\ 80 & 80 & 80 \end{smallmatrix}\big)$ and apply the vertical Sobel kernel
$G_y$ from Eq.~\eqref{ch32:sobel-x} by cross-correlation centered on the middle
pixel. Compute the response. Does its sign tell you whether the edge runs
dark-to-light going downward or upward? Now apply $G_x$; explain why you get
zero.
\end{exercise}

\begin{exercise}
\label{ch32:ex-sum-rule}
For each kernel below, sum the weights and state from the rule in
Section~\ref{ch32:blur-vs-diff} whether it blurs, differentiates, or sharpens:
(a) $\tfrac{1}{16}\big(\begin{smallmatrix} 1&2&1\\2&4&2\\1&2&1
\end{smallmatrix}\big)$; (b) $\big(\begin{smallmatrix} 0&-1&0\\-1&5&-1\\0&-1&0
\end{smallmatrix}\big)$; (c) the Laplacian of Eq.~\eqref{ch32:laplacian}.
Verify your answers by predicting what each does to a flat gray image.
\end{exercise}

\begin{exercise}
\label{ch32:ex-separable}
Show that the Gaussian kernel of Eq.~\eqref{ch32:gauss-kernel} is separable by
writing it as an outer product of two 1D Gaussians. Count the multiplies per
output pixel for a $7\times7$ Gaussian done as a single 2D convolution versus
two 1D convolutions. By what factor is the separable version cheaper?
\end{exercise}

\begin{exercise}
\label{ch32:ex-dog}
The \emph{difference of Gaussians} approximates the Laplacian-of-Gaussian:
$\text{DoG} = g_{\sigma_1} - g_{\sigma_2}$ with $\sigma_1 < \sigma_2$. Sketch
the 1D profile of $g_{\sigma_1}(x) - g_{\sigma_2}(x)$ and argue that it has the
Mexican-hat shape of Figure~\ref{ch32:fig-log}. Why is DoG so much cheaper to
compute than a true LoG in a multi-scale pyramid?
\end{exercise}

\begin{exercise}
\label{ch32:ex-noise}
Explain, using the frequency-domain view, why differentiating an image
amplifies noise and why pre-smoothing fixes it. Relate your answer to the
band-pass nature of the Sobel operator described in Section~\ref{ch32:sobel}.
\end{exercise}

\begin{exercise}
\label{ch32:ex-tensor}
Consider the structure tensor of Eq.~\eqref{ch32:structure-tensor} computed
over a window lying entirely on a straight vertical edge. Argue that one
eigenvalue is large and the other near zero. What does this imply for the
Harris response $R$? Why does this correctly say ``edge, not corner''?
\end{exercise}

\begin{exercise}
\label{ch32:ex-corr-vs-conv}
A library reports that it computes ``convolution,'' but feeding it the kernel
$(-1, 0, 1)$ and a step that rises left-to-right gives a \emph{negative} peak,
the opposite of what Eq.~\eqref{ch32:central-diff} predicts. Using the distinction
between convolution and cross-correlation in Section~\ref{ch32:2d-conv}, explain what happened and how to fix
your kernel so the signs come out as expected.
\end{exercise}

\begin{exercise}
\label{ch32:ex-cnn-params}
A convolutional layer with $64$ kernels of size $3\times3$ takes a single-
channel input. How many weights does it have (ignore biases)? Compare to a
fully connected layer mapping a $256\times256$ image to $64$ outputs. Use the
ratio to explain the ``weight sharing'' argument of Section~\ref{ch32:cnn}.
\end{exercise}

\begin{exercise}
\label{ch32:ex-imu-map}
You build a $6 \times 128$ ``channel $\times$ time'' array from one second of
accel+gyro data and convolve it with a $3\times3$ learned kernel. Describe in
words one spatiotemporal pattern such a kernel could learn to detect. Then
state plainly one thing this representation \emph{cannot} recover from the wrist
alone, citing the limits-of-recovery trap of Section~\ref{ch32:to-wrist}.
\end{exercise}

\begin{exercise}
\label{ch32:ex-learned-classic}
It is observed that the first layer of a CNN trained on photos learns filters
resembling oriented edge detectors. Give two independent reasons --- one about
the statistics of natural images, one about the architecture's built-in priors
--- for why this happens, drawing on the discussion of natural-image priors and learned filters in
Section~\ref{ch32:cnn}.
\end{exercise}
